% =====================================================================
%  Lời mở đầu và bảng đối chiếu với đề bài
% =====================================================================
\chapter*{Lời mở đầu}
\addcontentsline{toc}{chapter}{Lời mở đầu}
\markboth{Lời mở đầu}{Lời mở đầu}

Assignment 04 đã cài đặt phép tích chập bằng tay và so sánh ba khung phần mềm. Assignment 05 đặt một câu hỏi khác:
từ năm 1998 đến nay các kiến trúc CNN đã thêm những gì, và mỗi thứ thêm vào thực sự đổi được bao nhiêu? Để trả lời,
bài nhìn mạng nơ-ron tích chập như một hợp của các hàm, $f = f_L\circ\dots\circ f_1$. Với cách nhìn này, mỗi kiến
trúc mới chỉ là một cách thay một hàm, hoặc thay cách nối các hàm với nhau. Ví dụ, ResNet không đổi hàm tích chập;
nó thêm $x$ vào đầu ra của hai tầng tích chập.

Báo cáo đi theo đúng năm mục của đề. Chương 1 viết từng khái niệm (neuron, hàm kích hoạt, tích chập, pooling,
softmax, lan truyền ngược) thành một hàm NumPy và đối chiếu với PyTorch. Chương 2 đi qua chín kiến trúc từ LeNet đến ViT,
mỗi kiến trúc một lớp PyTorch ngắn. Chương 3 mô tả ba tập dữ liệu: CIFAR-10, CIFAR-100 và Diabetes Prediction. Chương 4
xây một CNN cơ bản và ba mô hình phát triển. Mỗi mô hình chỉ thêm đúng một cơ chế so với mô hình trước, rồi cả bốn được
huấn luyện trên cả ba tập. Chương 5 so sánh chúng theo chất lượng, chi phí, tốc độ hội tụ, lỗi và biểu diễn.

Kết quả không đi theo câu chuyện ``kiến trúc sau luôn tốt hơn kiến trúc trước''. Bước từ CNN cơ bản sang VGG tăng
\SL{compare/delta/cifar10/basic_vgg} điểm accuracy trên CIFAR-10 và \SL{compare/delta/cifar100/basic_vgg} điểm trên CIFAR-100.
Thêm đường tắt của ResNet không đổi gì trên CIFAR-10 và chỉ thêm \SL{compare/delta/cifar100/vgg_resnet} điểm trên CIFAR-100.
Thêm khối chú ý SE không cho lợi ích nhất quán: accuracy đổi \SL{compare/delta/cifar10/resnet_seresnet} điểm trên CIFAR-10
và \SL{compare/delta/cifar100/resnet_seresnet} điểm trên CIFAR-100. Trên dữ liệu bệnh tiểu đường, cả bốn mô hình có
ROC-AUC gần như bằng nhau, và mô hình rẻ nhất là lựa chọn hợp lý. Chương 5 giải thích vì sao: mỗi cơ chế chỉ có ích khi mô hình
gặp đúng vấn đề mà cơ chế đó được tạo ra để giải quyết.

Mọi con số trong báo cáo được sinh tự động từ tệp kết quả của notebook, mọi đoạn mã được trích thẳng từ mã nguồn, và
mọi hình được vẽ bằng TikZ từ dữ liệu, kể cả ảnh mẫu (vẽ lại từng điểm ảnh). Nếu chạy lại notebook, báo cáo sẽ đổi theo.
Phụ lục~\ref{ch:phuluc} ghi các lệnh để tái lập.

Em xin cảm ơn thầy \textbf{PGS.TS Trần Đình Quế} đã hướng dẫn học phần và ra đề bài này.

\vspace{8mm}
\begin{flushright}
\emph{Hà Nội, tháng 9 năm 2026}\par
\vspace{2mm}
\textbf{Nguyễn Duy Nghĩa}
\end{flushright}

\section*{Đối chiếu với đề bài}

\begin{table}[H]
  \centering
  \small
  \begin{tabular}{L{4.6cm}L{2.2cm}L{3.8cm}L{3.8cm}}
    \toprule
    Yêu cầu của đề & Chương & Notebook & Mã nguồn \\
    \midrule
    1. Khái niệm cơ bản của CNN qua hàm và hợp hàm, kèm code & \ref{ch:khainiem} & \texttt{00\_ham\_va\_hop\_ham} & \texttt{a05/concepts.py} \\
    2. Các mô hình phát triển của CNN, hàm và code & \ref{ch:kientruc} & \texttt{01\_kien\_truc\_phat\_trien} & \texttt{a05/blocks.py} \\
    3. Hai tập ảnh và một tập tiểu đường, có link và mô tả & \ref{ch:dulieu} & \texttt{02\_du\_lieu} & \texttt{a05/data.py} \\
    4. Code một CNN cơ bản và ba mô hình phát triển & \ref{ch:mohinh} & \texttt{03\_cifar10}, \texttt{04\_cifar100}, \texttt{05\_diabetes} & \texttt{a05/models.py}, \texttt{a05/train.py} \\
    5. So sánh và đánh giá & \ref{ch:sosanh} & \texttt{06\_so\_sanh} & \texttt{a05/metrics.py} \\
    \bottomrule
  \end{tabular}
\end{table}
